\documentclass[10pt,a4paper]{article}
\usepackage{amsmath,amssymb,amsthm,graphicx}
\usepackage{amsthm}
\usepackage{amsfonts, amsmath, amssymb, amsthm}
\usepackage{hyperref}
\usepackage{cleveref}
\def\TT{\mathbb{T}}
\def\PP{\mathbb{P}}
\def\DD{\mathbb{D}}
\def\QQ{\mathbb{Q}}
\def\WW{\mathbb{W}}
\def\MM{\mathbb{M}}
\def\HH{\mathbb{H}}

\newcommand{\Psf}{\mathsf{P}}
\newcommand{\Qsf}{\mathsf{Q}}
\newcommand{\Rsf}{\mathsf{R}}


\def\F{\mathcal{F}}
\def\U{\mathcal{U}}
\def\C{\mathcal{C}}
\def\I{\mathcal{I}}
\def\D{\mathcal{D}}
\def\L{\mathcal{L}}
\def\V{\mathcal{V}}
\def\A{\mathcal{A}}
\def\B{\mathcal{B}}
\def\P{\mathcal{P}}
\def\M{\mathcal{M}}
\def\Nc{\mathcal{N}}
\def\Q{\mathcal{Q}}

% Baire space, Cantor space, binary strings
\newcommand{\cs}{2^{\NN}}
\newcommand{\bs}{\NN^{\NN}}
\newcommand{\bstr}{2^{<\NN}}

\newcommand{\uh}{\upharpoonright}

% Turing machines
\newcommand{\converge}{\!\!\downarrow}
\newcommand{\diverge}{\!\!\uparrow}

% Forcing question
\newcommand{\qvdash}{\operatorname{{?}{\vdash}}}
\newcommand{\nqvdash}{\operatorname{{?}{\nvdash}}}

\newcommand{\ISig}{\mathsf{I}\Sigma^0}
\newcommand{\Robinson}{\mathsf{Q}}
\newcommand{\IPi}{\mathsf{I}\Pi^0}
\newcommand{\IDelta}{\mathsf{I}\Delta^0}
\newcommand{\Iopen}{\mathsf{I_{open}}}
\newcommand{\BSig}{\mathsf{B}\Sigma^0}
\newcommand{\BPi}{\mathsf{B}\Pi^0}
\newcommand{\BDelta}{\mathsf{B}\Delta^0}
\newcommand{\BCSig}{\mathsf{BC}\Sigma^0}
\newcommand{\BCPi}{\mathsf{BC}\Pi^0}
\newcommand{\BCDelta}{\mathsf{BC}\Delta^0}
\newcommand{\LSig}{\mathsf{L}\Sigma^0}
\newcommand{\LPi}{\mathsf{L}\Pi^0}
\newcommand{\CSig}{\mathsf{C}\Sigma}

\newcommand{\RCA}[0]{\mathsf{RCA}}
\newcommand{\WKL}[0]{\mathsf{WKL}}
\newcommand{\WWKL}[0]{\mathsf{WWKL}}
\newcommand{\ACA}[0]{\mathsf{ACA}}
\newcommand{\ATR}[0]{\mathsf{ATR}}
\newcommand{\PIOOCA}[0]{\Pi^1_1\mathsf{-CA}}
\newcommand{\CAC}[0]{\mathsf{CAC}}
\newcommand{\ADS}[0]{\mathsf{ADS}}
\newcommand{\RRT}[0]{\mathsf{RRT}}
\newcommand{\FS}[0]{\mathsf{FS}}
\newcommand{\PFS}[0]{\mathsf{PFS}}
\newcommand{\TS}[0]{\mathsf{TS}}
\newcommand{\GRRT}[0]{\mathsf{GRRT}}
\newcommand{\LRRT}[0]{\mathsf{LRRT}}
\newcommand{\EM}[0]{\mathsf{EM}}
\newcommand{\SEM}[0]{\mathsf{SEM}}
\newcommand{\HalfEM}[0]{\mathsf{Half}\mbox{-}\mathsf{EM}}
\newcommand{\RT}[0]{\mathsf{RT}}
\newcommand{\SRT}[0]{\mathsf{SRT}}
%\newcommand{\RRT}{\mathsf{RRT}}

\newcommand{\NN}[0]{\mathbb{N}}
\newcommand{\nat}[0]{\mathbb{N}}

\newcommand{\card}{\operatorname{card}}
\newcommand{\finsub}{\subseteq_{\mathtt{fin}}}

% reductions

\newcommand{\W}{\mathrm{W}}
\newcommand{\sW}{\mathrm{sW}}

\newcommand{\lex}{\mathrm{lex}}

% Barriers
\newcommand{\base}{\operatorname{base}}
\newcommand{\range}{\operatorname{ran}}

%typographically correct double quotes
\def\qt#1{``#1''}%
\newtheorem*{statement}{Statement}
\newtheorem{theorem}{Theorem}
\numberwithin{theorem}{section}
\newtheorem{maintheorem}[theorem]{Main Theorem}
\newtheorem{lemma}[theorem]{Lemma}
\newtheorem{question}[theorem]{Question}
\newtheorem{proposition}[theorem]{Proposition}
%\newtheorem{statement}[theorem]{Statement}
\newtheorem{remark}[theorem]{Remark}
\newtheorem{definition}[theorem]{Definition}
\newtheorem{corollary}[theorem]{Corollary}
\newtheorem{example}[theorem]{Example}

\newcommand{\R}{\mathcal{R}}
\renewcommand{\S}{\mathcal{S}}
\newcommand{\T}{\mathcal{T}}
\title{One computable Schreier coloring whose thin sets uniformly compute the omega Turing jump}
\author{Lorenzo Carlucci, Oriola Gjetaj, Quentin Le Houérou, Ludovic Levy Patey}
\date{arxiv:2412.11598v3; CC BY 4.0}
\begin{document}
\maketitle
\noindent\textit{Editorial note.} The full finite-jump coloring and oracle-removal lemmas are proof obligations, included before the final Schreier construction. Only the two explicitly stated classical computability tools are allowed without reproving them.\par\medskip
\noindent\textit{Editorial note.} A uniform jump modulus may be chosen nondecreasing by taking finite maxima; this makes the domination step at source696 explicit.\par\medskip
\noindent\textit{Editorial note.} The auxiliary n=1 finite coloring is necessarily constant0 and has no infinite thin set relative to its one-element codomain, so its asserted decoding property is vacuous. The meaningful finite-jump lemma683-697 has n>=2; the final decoding can choose x0>=4 as well as the printed inequalities, so it uses that proved range.\par\medskip
\noindent\textit{Editorial note.} The source726 phrase c is part of the range refers to the specified color codomain \{0,...,x0-k-1\}; the proof uses the inequality c<x0-k. It does not assume that each color is actually attained on the selected tail.\par\medskip
\noindent\textit{Editorial note.} The uniform decoder is given both H and an avoided color c, exactly as required by the counted theorem. No algorithm to find an avoided color from H alone is claimed.\par\medskip
\noindent\textit{Editorial scope.} One computable-instance lower bound for the large thin set theorem. All finite-jump coding lemmas and the uniform join argument are included. Reverse-mathematical equivalences, free-set reductions, rainbow cone avoidance and the commented older free-set proposition are not counted.\par\medskip
\noindent\textit{Allowed standard prerequisites (editorial).}\begin{itemize}
\item Uniform self-moduli for finite Turing jumps. For each positive k one can uniformly obtain an index for a nondecreasing function g\_k computable in the k-th jump of the empty set, and a decoding functional, such that any total function pointwise dominating g\_k computes that k-th jump through this functional. The indices are computable uniformly in k.. The precise classical modulus input invoked in source688; monotone closure of a modulus is permitted.. Provides only the starting modulus for each jump. The colorings, preservation under dimension increase, and one-color-avoidance decoding are proved in the solution.
\item Relativized Limit Lemma, uniformly in indices. Given an index of a total finite-valued X-prime-computable function f, one uniformly obtains total X-computable approximations f\_s with limit\_s f\_s(u)=f(u) for each finite input u.. Exact approximation form invoked in source701-708.. Used to replace the last finite tuple coordinate by an approximation stage, removing one jump while increasing arity by one.
\item Basic Turing reducibility and join coding. A computation uniformly obtaining the n-th jump for all n computes their countable join, the omega jump. Infinite sets are represented by characteristic functions and can uniformly enumerate their increasing elements.. Assemble the individual decoded jumps into one oracle computation from H together with an avoided color.
\end{itemize}\medskip
There exist four other subsystems which, together with $\RCA_0$, calibrate the strength of most theorems. These subsystems are known as the \qt{Big Five} (see Mont\'alban~\cite{montalban_open_2011}). Among these, $\ACA_0$ (Arithmetic Comprehension Axiom) extends the axioms of $\RCA_0$ with the comprehension scheme for all arithmetic formulas. Based on the correspondence between computability and definability, it is equivalent to stating the existence of the Turing jump of any set ($\forall X \exists Y (Y = X')$).

We shall also consider two lesser-known stronger variants of $\ACA_0$, namely, $\ACA_0'$ and $\ACA_0^+$. The system $\ACA_0'$ extends $\RCA_0$ with the axiom stating closure under all finite jumps; the system $\ACA_0^+$ extends $\RCA_0$ with the axiom stating the existence of the $\omega$-jump of any set. The $\omega$-jump of a set~$X$ is the set $X^{(\omega)} = \bigoplus_n X^{(n)}$. The system $\ACA_0^+$ is famous for being the best known upper bound to Hindman's theorem.%\lorenzo{I added mention of ACA0' since we mention it at least once.}

Let us first fix some notation and recall the definition of these principles for colorings of finite subsets of fixed size. We denote by $\NN$ the set of natural numbers and by $\NN^+$ the set of positive natural numbers. For $X\subseteq \NN$ and $n \in \NN$
we denote by $[X]^{n}$ the set of all $n$-subsets of $X$. We denote by $[X]^{<\omega}$ the set of all finite subsets of $X$. We always assume that sets are presented in increasing order and make no distinction between a set and the sequence of its elements in increasing order. We write $x \cdot y$ for the concatenation of two sequences $x$ and $y$. We identify a natural number with the set of its predecessors, so that, if $k\geq 1$ we can write $f: X \to k$ to declare a function from $X$ to $\{0, 1, \dots, k-1\}$. A function of this type is often called a coloring of $X$ in $k$ colors.

\begin{definition}[(Exactly) $\omega$-large sets] 
A finite $s\subseteq \NN$ is $\omega$-large if $|s| \geq 1+\min s$ and is exactly $\omega$-large if $|s| = 1+\min s$.
\end{definition}

The $\omega$-large sets are also known as {\em relatively large} sets in the literature. For an infinite $X\subseteq \NN$ we denote by $[X]^{!\omega}$ the family of all exactly $\omega$-large subsets of $X$.

\begin{definition}[Large Thin Set Theorem]\label{thm:thin}
For every coloring $f:[\NN]^{!\omega} \to \NN$, there exists an infinite set $H\subseteq \NN$ such that $f([H]^{!\omega}) \neq \NN$. The set $H$ is called \emph{thin} for $f$. We abbreviate this statement by $\TS^{!\omega}$.
\end{definition}

\begin{definition}
Let $C$ be any non-empty set. Given a coloring $f : [\NN]^n \to C$, a set~$H \subseteq \NN$ is \emph{$f$-thin for color~$c \in C$} if $c \not \in f([H]^n)$. A set ~$H \subseteq \NN$ is $f$-thin (or thin for $f$) if $H$ is $f$-thin for some color $c\in C$.
\end{definition}

The definition of $f$-thin set depends on the choice of codomain of the function. It will always be clear from the context. 

\begin{theorem}\label[theorem]{ts-omega-computes-omega-jump}
    There exists a computable coloring $f : [\NN]^{!\omega} \to \NN$ such that every infinite $f$-thin set~$H$ computes $\emptyset^{(\omega)}$. Moreover, this computation is uniform in~$H$ and an avoided color~$c$.
\end{theorem}

\begin{lemma}\label{lem:arrays}
    There exists two computable arrays $(e_{n,k})_{n,k \in \NN}$ and $(d_{n,k})_{n,k \in \NN}$ of Turing indexes such that for every $n, k \in \NN^+$ with $n \geq 2$, $\Phi_{e_{n,k}}^{\emptyset^{(k)}}$ is a coloring $f_n^k : [\NN]^{n} \to n$ such that for every infinite $f_n^k$-thin set $H$, $\Phi_{d_{n,k}}^{H} = \emptyset^{(k)}$.
\end{lemma}

\begin{proof}
Let $k\in\NN^+$. Let $g_k$ be a uniform modulus of the set $\emptyset^{(k)}$. Note that $\Delta^0_{k+1}$-indexes for each function $g_k$ can be found computably, uniformly in~$k$. 

For $n \geq 2$ let $f_n^k(x_0, \dots,x_{n-1}) = n - 1$ if $g_k(x_0) \leq x_1$ and if $g_k(x_0) > x_1$, let $f_n^k(x_0, \dots,x_{n-1}) = i$ for $i < n - 1$ the largest value such that $g_k(x_0) > x_{i+1}$.

Let $H$ be an infinite $f_n^k$-thin set for some color $c$. The color $c$ cannot be $n-1$ as for a given $x_0 \in H$ there exists $x_1 <\dots < x_{n-1} \in H \setminus \{0, \dots, g_k(x_0)\}$, hence $f^k_n(x_0, \dots, x_{n-1}) = n-1$.

We claim $H$ is thin for color $n-2$. Indeed, assume by contradiction that there exists some tuple $x_0 < \dots < x_{n-1} \in H$ such that $g_k(x_0) > x_{n-1}$ so that $f^k_n$ takes color $n-2$ on $\{x_0, \dots, x_{n-1}\}$. Since $c < n-2$ and $H$ is infinite, there exists some $y_{c+2}< \dots < y_{n-1} \in H \setminus \{0, \dots ,g_k(x_0)\}$, therefore $f^k_n(x_0, \dots, x_{c+1},y_{c+2}, \dots, y_{n-1}) = c$, contradicting our assumption that $H$ is $f_n^k$-thin for $c$.

Write $H = \{x_0 < x_1 < \dots \}$. For every $i \in \NN$, $g_k(x_i) \leq x_{i + n - 1}$, hence $H$ computes a function dominating $g_k$ and therefore computes $\emptyset^{(k)}$. This computation can be done uniformly in the set $H$, $k$ and $n$. 
\end{proof}

\begin{lemma}
    There exists a functional $\Gamma$ such that for every set $X$ and every index $e$, if $\Phi_e^{X'}$ is a coloring $f : [\NN]^n \to \ell$, then $\Gamma^X(e)$ is a coloring $g : [\NN]^{n+1} \to \ell$ such that if an infinite set $H$ is $g$-thin for a color $c$, then $H$ is $f$-thin for the same color $c$. 
\end{lemma}

\begin{proof}
Consider a set $X$ and an index $e$ such that $\Phi_e^{X'}$ is a coloring $f : [\NN]^n \to \ell$ for some $n, \ell \in \NN^+$. Let $(f_s)_{s\in \NN}$ be a $\Delta_2^0(X)$-approximation of $f$ (indexes for such an approximation, where the $f_s$ are seen as $X$-computable functions, can computably be found uniformly in $X$ and $e$). Finally, consider the following $X$-computable coloring $g$ defined by $g(x_0, \dots, x_{n}) = f_{x_n}(x_0, \dots, x_{n-1})$ (again, the construction is uniform).

Let $H \subseteq \NN$ be an infinite set. If for some $x_0 < \dots < x_{n-1} \in H$ we have $f(x_0, \dots, x_{n-1}) = c$ for some color $c$, then $f_s(x_0, \dots, x_{n-1}) = c$ for every $s$ bigger than a certain threshold. Thus, as $H$ is infinite, there exists some $x_n \in H$ such that $g(x_0, \dots, x_n) = c$.  
\end{proof}

Combining these two lemmas, we get the following, where the array $(d_{n,k})_{n,k\in\NN}$ is as in~\Cref{lem:arrays}.

\begin{lemma}\label[lemma]{lem:uniform-computable-thin-coloring}
    There exists a uniformly computable sequence of colorings $f_{n,k} : [\NN]^{n+k} \to n$, for $n, k \in \NN^+$, such that for every infinite $f_{n,k}$-thin set $H$, $\Phi_{d_{n,k}}^H = \emptyset^{(k)}$.
\end{lemma}

\begin{proof}
Let $f : [\NN]^{!\omega} \to \NN$ be defined as follows (where the functions $f_{n,k}$ are from \Cref{lem:uniform-computable-thin-coloring}): for $x_0 < \dots < x_{x_0}$ in $\NN$ with $x_0 > 1$, let $f(x_0, \dots, x_{x_0})= f_{x_0 - k, k}(x_1, \dots, x_{x_0})$ for $k = \lceil \frac{x_0}{2} \rceil$; for $s \in [\NN]^{!\omega}$ wit $\min s \leq 1$ set $f(s) = 0$. By \Cref{lem:uniform-computable-thin-coloring}, $f$ is indeed computable.

Let $H$ be an infinite $f$-thin set for some color $c$ and let $n \in \NN$. Since $H$ is infinite, there exists some $x_0 \in H$ such that, by letting $k = \lceil \frac{x_0}{2} \rceil$, we have $x_0 - k > c$ and $k \geq n$. The set $G = H \setminus \{0, \dots, x_0\}$ is infinite and $f$-thin for $c$. By definition of $f$, $G$ is $f_{x_0 - k, k}$-thin for the color~$c$. Note that $c$ is part of the range of $f_{x_0 - k, k}$ as $x_0 - k > c$. By \Cref{lem:uniform-computable-thin-coloring}, since $k \geq n$, $G \geq_T \emptyset^{(n)}$. Since $G$ is obviously $H$-computable, we have $H\geq_T \emptyset^{(n)}$.
This computation can be done uniformly in~$H$, $c$ and $n$, thus $H \geq_T \emptyset^{(\omega)}$ uniformly in $H$ and~$c$.
\end{proof}

\begin{thebibliography}{1}
\bibitem{montalban_open_2011} Antonio Montalb\'an, Open questions in reverse mathematics,
\emph{Bulletin of Symbolic Logic} 17(03) (2011), 431--454.
\end{thebibliography}
\end{document}
